\section{Dalvik e ART: registradores, AOT e JIT}
\label{sec:hist-dalvik}

\paragraph{Dalvik: uma VM de registradores para Java.}
O Android usa Java como linguagem, mas não roda uma JVM: o runtime
executa bytecode no formato DEX (\citetitle{doc:android-art}).
A especificação diz: \enquote{The machine is
register-based, and frames are fixed in size upon creation}; cada frame
tem um número fixo de registradores de 32~bits, valores de 64~bits usam
pares, e o desenho tenta \enquote{approximately imitate common real
architectures and C-style calling conventions}
(\citetitle{doc:android-dalvik-bytecode}).

\paragraph{Pilha ou registradores?}
\textcite{paper:vm-showdown} mediram a troca que o Dalvik fez. Eles
traduziram bytecode JVM de pilha para um formato de registradores dentro
de uma JVM interpretada (Sun J2ME CDC~1.0), com SPECjvm98 e Java Grande.
Resultado, medido pelos autores (Trinity College Dublin e TU Wien, 2005):
a versão de registradores executa, em média, mais de 47\% menos
instruções da VM, com bytecode cerca de 25\% maior; num Pentium~4, leva
em média 32,3\% menos tempo com despacho por \texttt{switch} em C e cerca
de 26,5\% menos com \emph{threaded dispatch}. Note o que foi medido: um
\emph{interpretador}. Num JIT, o bytecode é só a entrada do compilador,
que gera código para os registradores da máquina real; o C2, por
exemplo, faz alocação de registradores por coloração de grafos
\parencite{paper:hotspot-server-compiler}. A vantagem medida vale para a
fase interpretada, não para o código quente.

\paragraph{ART: AOT na instalação, JIT com perfil.}
O sucessor do Dalvik é o ART. A documentação diz que \enquote{ART and
Dalvik are compatible runtimes running DEX bytecode} e que, na
instalação, o ART compila os apps com a ferramenta \texttt{dex2oat}; o GC
é descrito como majoritariamente concorrente, com uma só pausa
(\citetitle{doc:android-art}). O ART também tem um JIT
\enquote{with code profiling}: um método pode estar interpretado,
compilado pelo JIT ou compilado AOT, e um daemon usa o perfil coletado
pelo JIT para guiar a compilação AOT, o que a documentação chama de
\emph{profile-guided compilation} (\citetitle{doc:android-art-jit}).

\paragraph{O que isso explica na JVM.}
\begin{itemize}
\item Que o bytecode de pilha da JVM é um formato de distribuição, não
  de execução: Dalvik e ART pegam o mesmo código Java e o executam em
  outra forma.
\item Que AOT e JIT não são rivais. O ART usa o perfil do JIT para
  compilar antes; no OpenJDK, o Projeto Leyden vai na mesma direção com o
  cache AOT da seção~\ref{sec:jvm-vs-nativo} e o perfil de métodos
  antecipado do \citetitle{doc:jep515}.
\end{itemize}
